% P2 related-work paragraph (draft for the preregistration; checked by citation-verifier; bib: reports/p2/related_work.bib)
Causal tracing and activation patching are used to locate where factual associations are processed, but where a fact is localized need not be where it is best edited \cite{hase2023}.
Chen et al.\ argue that the knowledge-localization assumption does not hold for all facts and propose query localization instead \cite{chen2024}.
Hochman et al.\ find factual retrieval to be redundant, distributed and non-contiguous across layers, and report significant rank discordance between models in per-relation path length \cite{hochman2026}.
Gupta et al.\ report that the first token of multi-token answers emerges later in depth than single-token answers \cite{gupta2025}.
Agarwal finds that relation information becomes causally active before entity information during factual recall \cite{agarwal2026}.
Rohweder et al.\ predict how linearly decodable a relation is from the subject representation alone, without the relation \cite{rohweder2026}.
Elhelo et al.\ find, with behavioural and parameter-localization analyses, that language models encode a fact's knowledge in a task-specific manner, so that different tasks on the same fact rely on different parameters \cite{elhelo2026}.
Benz et al.\ present KnowledgeDebugger, a tool for exploring knowledge localization and editing in transformers \cite{benz2026}.
To our knowledge, none of these estimates how much of a fact's causal resolution depth is a property of the fact and how much of the prompt wording; the closest are Elhelo et al.\ (task dependence of a fact's parameters) and Agarwal (eight prompt families).
